\subsection{Haefliger's modular-group coinvariant equation: a complete continuous result}
\label{dev:1981-haefliger}

The question in 1981-555.1 is exact smooth solvability on the projective circle
\(X=\mathbb {RP}^{1}=\mathbb R\cup\{\infty\}\): for every real
\(f\in C^\infty(X)\), do there exist real \(u,v\in C^\infty(X)\) such that
\[
 f=u-u\circ T+v-v\circ S,
 \qquad T(x)=x+1,\qquad S(x)=-1/x?
\]
Both maps extend smoothly across the apparent affine-coordinate singularities.
The notebook calls exact solvability unknown and asserts density of the image
in the \(C^\infty\) topology. Those are two separate assertions. In particular,
density in a Fr\'echet space does not imply surjectivity.

\paragraph{Literature and the scope of this development.}
The search documented in the research log located Haefliger's 1980 paper on
minimal leaves, and the later invariant-functional framework of
\citet[abstract]{abouqatebelkacimi2000invariant}. The latter's
compact-group result does not cover this noncompact discrete group action.
Neither lead supplied a verified theorem answering the displayed smooth
equation. The subject connection with minimal leaves is recorded in the notes;
no historical explanation of the question's subsequent neglect is inferred.
The following result is a bounded contribution proved here, with no claim of
historical priority. It solves the exact equation on continuous functions. It
does not prove a regularity theorem for its solutions.

\begin{lemma}[No invariant probability measure]\ours
There is no Borel probability measure on \(X\) invariant under both \(T\)
and \(S\).
\end{lemma}
\begin{proof}
If \(\nu\) is \(T\)-invariant, the pairwise disjoint Borel sets
\([k,k+1)\), \(k\in\mathbb Z\), have the same measure. Finiteness forces
each measure to be zero, and countable additivity then gives
\(\nu(\mathbb R)=0\). Thus \(\nu=\delta_\infty\). Since
\(S(\infty)=0\), this probability is not \(S\)-invariant.
\end{proof}

\begin{lemma}[A uniform defect of invariance]\ours
Let \(\mathcal M(X)\) be the real finite signed regular Borel measures, with
total-variation norm \(\|\cdot\|_{\rm TV}\), and let \(g_*\mu\) denote
pushforward. There is a number \(\epsilon>0\) such that every
\(\mu\in\mathcal M(X)\) satisfies
\[
 \|\mu-T_*\mu\|_{\rm TV}+\|\mu-S_*\mu\|_{\rm TV}
 \ \geq\ \epsilon\|\mu\|_{\rm TV}.
 \tag{H1}
\]
\end{lemma}
\begin{proof}
The probabilities on compact \(X\) form a weak-* compact subset of
\(C(X)^*\). Pushforward by a homeomorphism is weak-* continuous and the
total-variation norm is weak-* lower semicontinuous. Consequently
\[
 \epsilon:=\inf_{\nu\ \mathrm{probability}}
 \bigl(\|\nu-T_*\nu\|_{\rm TV}+\|\nu-S_*\nu\|_{\rm TV}\bigr)
\]
is attained. If it were zero, a minimizer would contradict the preceding
lemma. Thus \(\epsilon>0\).

For signed measures \(\alpha,\beta\),
\(\||\alpha|-|\beta|\|_{\rm TV}\leq\|\alpha-\beta\|_{\rm TV}\).
Indeed, taking Radon--Nikodym derivatives with respect to
\(|\alpha|+|\beta|\) reduces this assertion to
\(\bigl||a|-|b|\bigr|\leq|a-b|\) and integration. A homeomorphism also
satisfies \(|g_*\mu|=g_*|\mu|\). Apply the positive-measure estimate to
\(|\mu|/\|\mu\|_{\rm TV}\), when \(\mu\ne0\), and use these two facts
for \(g=T,S\). The zero measure causes no exception.
\end{proof}

\begin{theorem}[Exact solvability for continuous functions]\ours
For every real \(f\in C(X)\), there are real \(u,v\in C(X)\) with
\[
 f=u-u\circ T+v-v\circ S.
 \tag{H2}
\]
If \(0<\epsilon'<\epsilon\), the solutions can be chosen to satisfy
\[
 \max(\|u\|_\infty,\|v\|_\infty)
 \leq \frac{2}{\epsilon'}\|f\|_\infty.
 \tag{H3}
\]
Here \(\epsilon\) is the positive constant in (H1); the theorem does not
claim an explicit numerical value for it or a linear choice of solutions.
\end{theorem}
\begin{proof}
Equip \(E=C(X)\oplus C(X)\) with the maximum norm and define the bounded
operator \(D:E\to C(X)\) by the right side of (H2). Under the Riesz
identification \(C(X)^*=\mathcal M(X)\), its adjoint is
\[
 D^*\mu=(\mu-T_*\mu,\mu-S_*\mu).
\]
The dual norm on \(E^*\) is the sum of the two total-variation norms.
Thus (H1) states precisely that \(D^*\) is bounded below by \(\epsilon\).

Let \(B_E\) be the closed unit ball and let
\(K=\overline{D(B_E)}\), with norm closure. For every \(\mu\),
\[
 \sup_{h\in K}\langle\mu,h\rangle
 =\|D^*\mu\|\geq\epsilon\|\mu\|_{\rm TV}.
\]
The Hahn--Banach separation theorem for closed convex sets implies
\(\epsilon' B_{C(X)}\subset K\) for every \(0<\epsilon'<\epsilon\).
To see the strict-inequality point explicitly, a hypothetical separating
functional for an element of \(\epsilon' B_{C(X)}\setminus K\) would
have supremum on \(K\) at least \(\epsilon\|\mu\|\), while its value
on that element is at most \(\epsilon'\|\mu\|\), a contradiction.

Set \(r_0=f\). Given \(r_j\ne0\), the inclusion just proved permits
a pair \(z_j\in E\) with
\[
 \|z_j\|_E\leq\|r_j\|_\infty/\epsilon',\qquad
 \|r_j-Dz_j\|_\infty\leq\tfrac12\|r_j\|_\infty.
\]
Put \(r_{j+1}=r_j-Dz_j\); if a residual is zero, take all subsequent
pairs zero. Then \(\|r_j\|_\infty\leq2^{-j}\|f\|_\infty\) and
\(\sum_j\|z_j\|_E\leq2\|f\|_\infty/\epsilon'\). Completeness gives
\(z=\sum_j z_j=(u,v)\in E\), and continuity of \(D\) gives
\(Dz=f\). This also proves (H3).
\end{proof}

\begin{proposition}[An exact reduction of the smooth question]\ours
For \(C_gw=w\circ g\), put \(A=(I+C_S)/2\). A smooth function \(f\)
has a smooth decomposition (H2) if and only if
\[
 Af\in\operatorname{ran}\bigl(A(I-C_T):C^\infty(X)\to
 C^\infty(X)^S\bigr),
 \tag{H4}
\]
where \(C^\infty(X)^S\) denotes the \(S\)-invariant functions.
\end{proposition}
\begin{proof}
Since \(S^2=I\) on \(\mathbb {RP}^{1}\), \(A\) is a projection and
\(A(I-C_S)=0\). Applying it to (H2) proves necessity. Conversely, choose
\(u\) with \(A(I-C_T)u=Af\). The remainder
\(r=f-(I-C_T)u\) satisfies \(Ar=0\), hence \(C_Sr=-r\).
Putting \(v=r/2\) gives \((I-C_S)v=r\) and proves sufficiency.
\end{proof}

\paragraph{Second check and remaining obstruction.}
An independent audit re-derived the measure argument and checked all norms.
It also gave a separate functional-analytic check: (H1) makes
\(\operatorname{ran}D^*\) norm closed and \(\ker D^*=0\); the Banach
closed-range theorem makes \(\operatorname{ran}D\) closed, while the
annihilator identity makes it dense, so \(D\) is onto. This agrees with
the constructive Hahn--Banach argument above. The use of total variation
of a signed measure is essential: mere absence of invariant signed
measures would show density, not a quantitative lower bound or
surjectivity.

The proof controls only uniform norms. Uniformly convergent series of
continuous solutions need not have bounded derivatives, and smoothing
does not commute with the two nonlinear coordinate actions in a way that
this argument estimates. In the smooth setting the dual obstructions are
distributions, not only measures. The notebook's claim of smooth density,
even if independently established, would still leave a closed-range or
regularity estimate to prove. Thus the strongest verified result in this
development is exact \(C^0\) solvability; the strongest formulation of the
original problem remains exact \(C^\infty\) solvability, retained as an
unresolved question rather than promoted to a theorem.

One further bounded investigation can target (H4): establish a tame
estimate for the averaged parabolic difference operator near the fixed
point \(\infty\), with a finite and explicitly controlled derivative
loss, or construct an obstruction to such an estimate. That is a precise
analytic task. The present argument supplies neither outcome, and the
absence of a verified resolution in this search is not evidence that no
resolution exists.

\subsection{Independent finite checks of two source emendations}

These checks support annotations 1981-548.1--2 and 1981-546.15; they are not
additional open problems. The partition theorem itself is the known identity
in \citet[Chapter I, Section 1, (1.7)--(1.7')]{macdonald1995symmetric}.

\begin{example}[Exact partition enumeration]\oursComputed
The supplied Python program enumerates partitions in every rectangle
\(0\leq m,n\leq6\) in two ways: weakly decreasing row lengths, and
\(n\)-element subsets of \(\{0,\ldots,m+n-1\}\), inverted as beta numbers.
Both enumerations agree and contain \(\binom{m+n}{n}\) partitions for each
rectangle. For all 3431 partitions across these 49 rectangles, the two
complementary exponent sets are disjoint and fill that interval. A separate
run on \(0\leq m,n\leq4\), with 251 partitions, agrees. The negative control
\(m=n=1\), \(\lambda=(1)\), makes the written generating-function left side
\(t+t^2\), whereas the required right side is \(1+t\). The corrected inverse
argument and conjugation give \(1+t\). These finite checks detect the source
error; they are not the proof for arbitrary rectangles.
\end{example}

\begin{proposition}[The simplex edge-length Jacobian]\ours
Order the independent Gram entries by diagonals first and then off-diagonals.
Write \(s_{ij}=\ell_{ij}^2\) for all \(n(n+1)/2\) edges of a labelled
\(n\)-simplex, with vertex zero distinguished. The absolute Jacobians are
\[
 \left|\det\frac{\partial(g_{ij})_{1\leq i\leq j\leq n}}
                         {\partial(s_{ij})_{0\leq i<j\leq n}}\right|
 =2^{-n(n-1)/2},
 \qquad
 \left|\det\frac{\partial(g_{ij})}{\partial(\ell_{ij})}\right|
 =2^n\prod_{0\leq i<j\leq n}\ell_{ij}.
\]
Consequently the density in equation (23) needs the product of lengths if
its differentials are length differentials.
\end{proposition}
\begin{proof}
The diagonal formula is \(g_{ii}=s_{0i}\). Each off-diagonal coordinate is
\(g_{ij}=(s_{0i}+s_{0j}-s_{ij})/2\). In the stated ordering the Jacobian is
block triangular, with diagonal blocks the identity and
\(-\tfrac12 I_{n(n-1)/2}\). This proves the first identity. The change from
lengths to squared lengths has diagonal Jacobian entries \(2\ell_{ij}\).
Their product gives the second identity. Finally
\(\det g=(n!\operatorname{vol}\sigma)^2\), so multiplying the invariant
matrix density by this Jacobian gives, up to a fixed constant,
\(\operatorname{vol}(\sigma)^{-(n+1)}\prod\ell_{ij}\,d\ell_{ij}\).
\end{proof}

\begin{example}[Second check of the Jacobian]\oursComputed
Independent symbolic differentiation of the full coordinate map in dimensions
\(n=1,2,3,4,5\) gives signed squared-length Jacobians
\(1,-1/2,-1/8,1/64,1/1024\), respectively. Division of the absolute
length-coordinate determinant by the product of all edge lengths gives
\(2,4,8,16,32\). An independent positive-definite test uses based orthonormal
edge vectors, whose pairwise squared distances are one or two, and recovers
the identity Gram matrix exactly in every dimension. For a single edge,
\(g=\ell^2\) gives \(g^{-1}dg=2\,d\ell/\ell\), while the written formula
would give \(d\ell/\ell^2\). The complete code and its output are included.
\end{example}
